\section{Axioms, Semantics, And Public Data}\label{sec:mathematical-framework}

This section fixes the semantic and axiomatic surfaces used by the problem.
It comes before the public artifact and commitment layer because those objects
are only meaningful after the public/private boundary is fixed.

\subsection{Fibrational Semantics}

The semantic layer is motivated by Helfer's first-order homotopical logic
\cite{helfer2019fol}. In this paper, the corresponding Lean files are used as
a formalized interface for the semantic structure needed by \HE{}, not as a
claim that every external theorem from Helfer's paper has been reproved.

\subsection{Fibrational Identity Objects}

Public equality is not strict equality of hidden representatives. It is mediated
by public classifiers, path data, and equivalence of the observable structure.
This is the semantic source of the hidden-operator model used in HCP.

\subsection{Homotopy Classes}

Public shadows may identify multiple hidden representatives while preserving
the invariants required by a verifier. At this stage the shadow is only a
public classifier or fingerprint; finite-code projection is introduced later.

\subsection{Axiom Surfaces}

The Lean development separates assumptions and theorem surfaces into small
modules. The relevant early interfaces are:
\begin{enumerate}[label=(\roman*)]
  \item foundational cochain and Hodge assumptions;
  \item fibration and homotopy-class interfaces;
  \item Heyting/oracle semantics for public judgments;
  \item sheaf-coherence and spectral interfaces;
  \item cochain-commitment payload and security assumptions.
\end{enumerate}

\subsection{Exact Structural Invariants}

The exact gates used by the Lean development are not a separate source of
structure. They package the same foundational and sheaf-coherence hypotheses
at the point where an object is admitted. In Lean, \texttt{ExactGateBundle} is
the foundation interface and \texttt{SheafExactGateBundle} is the
sheaf-coherence interface. The theorem
\texttt{structural\_preconditions\_of\_sheaf\_exact\_gates} extracts the two
structural consequences used later: the cochain differential squares to zero
and the Hodge Laplacian is the one induced by the coboundary data.

\begin{proposition}[Exact structural preconditions]\label{prop:exact-structural-preconditions}
  An accepted sheaf-coherence surface supplies
  \[
    d^2=0
    \qquad\text{and}\qquad
    \Delta=d^\ast d+dd^\ast
  \]
  for the cochain complex under discussion.
\end{proposition}

\subsection{Heyting And Oracle Judgment Surface}

The public judgment layer is also fixed here, because HCP uses public
classifiers and oracle-visible predicates before any finite-code projection is
introduced. Lean separates an \texttt{OracleQuery} from transaction cochain
commitments: its \texttt{id\_commit} is an identity hash obtained from the
public Hodge projection and quantization path, not a cochain commitment
payload.

The oracle predicate checks public compatibility data such as the vertex
identifier, Betti number, identity hash, and Tsirelson-style spectral judgment.
The induced judgment is then mapped into the three-valued Heyting surface used
by the acceptance predicate. The stability axiom \texttt{OracleEta2} records
that equal public identity inputs produce the same identity hash; it is a
determinism condition for public identity observations, not a uniqueness or
binding theorem for cochain commitments.

This modular split is part of the mathematical statement. It prevents runtime
facts, cryptographic assumptions, public classifiers, and Hodge-theoretic
claims from being conflated. Exact and oracle invariants are used here as
mathematical acceptance surfaces. Consensus composition and chain/runtime
enforcement are recorded only at the formalization boundary; their full
CS/protocol treatment is deferred to the sequel.
